\section*{Latihan dan Refleksi}
\addcontentsline{toc}{section}{Latihan dan Refleksi}

\subsection*{Latihan}
\begin{enumerate}[label=\arabic*.,leftmargin=*]
  \item Jelaskan perbedaan assembler, compiler, interpreter, dan
        cross-compiler. Berikan satu contoh untuk masing-masing.
  \item Sebutkan enam fase kompilasi dan jelaskan masukan serta
        keluaran setiap fase.
  \item Mengapa tabel simbol disebut pusat informasi kompilator?
  \item Berikan tiga contoh kesalahan leksikal, tiga kesalahan
        sintaksis, dan tiga kesalahan semantik.
  \item Diberikan program C sederhana, tuliskan aliran token beserta
        lokasi setiap token.
  \item Gambarkan pohon sintaks untuk program tersebut.
  \item Tuliskan kode antara tiga alamat untuk program tersebut.
  \item Lengkapi tabel penelusuran fase dan jelaskan satu sel yang
        paling sulit Anda isi.
\end{enumerate}

\subsection*{Refleksi}
\begin{reflectionbox}
\begin{enumerate}[label=\arabic*.,leftmargin=*]
  \item Setelah membaca bab ini, fase mana yang paling menarik bagi
        Anda? Mengapa?
  \item Apakah Anda dapat menjelaskan mengapa compiler dan interpreter
        dapat menghasilkan perilaku yang berbeda untuk program yang
        sama?
  \item Apa satu hal yang akan Anda lakukan berbeda saat membaca
        program setelah memahami token, pohon sintaks, dan kode antara?
\end{enumerate}
\end{reflectionbox}
